% ============================================================
% Background & Related Work.
% Background: only what a reviewer needs to follow the paper —
% not a textbook chapter. Related work: organize by THEME, not
% chronologically, and always end a group with how your work
% differs. See references/writing-guide.md.
% ============================================================
\section{Background \& Related Work}
\label{sec:background}

\subsection{Background}
\label{subsec:background}

% 2–4 concepts the reader must know, each in one paragraph:
% define it, cite its origin, and say where it is used in THIS
% paper. If a concept needs a formal definition, state it here
% and give it a numbered equation so it can be referenced:
% \begin{equation}
%     \label{eq:problem}
%     \min_{x \in \mathcal{X}} f(x) \quad \text{s.t.} \quad g(x) \le 0
% \end{equation}
Define the concepts and notation the rest of the paper relies on.

\subsection{Related Work}
\label{subsec:related}

% Group related work by approach/theme (e.g., "graph-based
% methods", "learning-based methods"), one paragraph per group:
%   1. One sentence: what unites this group of work.
%   2. 2–4 sentences: key works, cited as \cite{a,b}, with the
%      specific technique each contributes.
%   3. One sentence: the shared limitation YOUR work addresses.
% Never dump citations without saying what they did. A
% comparison table helps when there are 3+ close works:
%
% \begin{table}[htbp]
%     \caption{Comparison with closely related work}
%     \label{tab:related-comparison}
%     \centering
%     \begin{tabular}{lccc}
%         \toprule
%         Approach & Property A & Property B & Scales \\
%         \midrule
%         Method X \cite{example2023} & \checkmark & $\times$ & $\times$ \\
%         Method Y \cite{conference2022} & $\times$ & \checkmark & $\times$ \\
%         \textbf{Ours} & \checkmark & \checkmark & \checkmark \\
%         \bottomrule
%     \end{tabular}
% \end{table}

Discuss how existing work addresses the problem and where this work is positioned.
